\usepackage[a4paper,margin=21mm,top=24mm,bottom=24mm]{geometry}
\usepackage{fontspec}
\usepackage{microtype}
\usepackage{graphicx}
\usepackage{xcolor}
\usepackage{booktabs}
\usepackage{array}
\usepackage{longtable}
\usepackage{caption}
\usepackage{enumitem}
\usepackage{fancyhdr}
\usepackage{amssymb}
\usepackage{amsmath}
\usepackage[hidelinks]{hyperref}

% STIX Two Text ships on this machine as a variable font whose Bold named
% instance XeTeX cannot select, so the bold shape has to be synthesised.
% Without this the abstract, which is set bold, renders at regular weight.
\setmainfont{STIX Two Text}[Numbers=Proportional,AutoFakeBold=1.6]
\newfontfamily\sffont{Helvetica Neue}
\newfontfamily\monofont{Menlo}[Scale=0.86]

\definecolor{ink}{HTML}{111111}
\definecolor{muted}{HTML}{6f6f6f}
\definecolor{rule}{HTML}{c9c9c9}
\definecolor{navy}{HTML}{2b5d8a}
\color{ink}

\setlength{\parindent}{1.4em}
\setlength{\parskip}{0pt}
\linespread{1.06}

\makeatletter
\renewcommand\section{\@startsection{section}{1}{\z@}{13pt}{4pt}%
  {\sffont\bfseries\fontsize{10.5}{12}\selectfont}}
\renewcommand\subsection{\@startsection{subsection}{2}{\z@}{9pt}{2pt}%
  {\sffont\bfseries\fontsize{9}{11}\selectfont\color{ink}}}
\renewcommand\subsubsection{\@startsection{subsubsection}{3}{\z@}{6pt}{-0.6em}%
  {\sffont\bfseries\fontsize{8.6}{10}\selectfont}}
\makeatother

\DeclareCaptionFont{capsf}{\sffont}
\captionsetup{font={footnotesize},labelfont={capsf,bf},labelsep=none,
  justification=justified,singlelinecheck=false,skip=5pt}
\captionsetup[figure]{name={Fig.}}
\captionsetup[table]{name={Table}}

\newcommand{\gene}[1]{\textit{#1}}
\newcommand{\pep}[1]{{\monofont #1}}
\newcommand{\hla}[1]{{\monofont #1}}
\newcommand{\figlead}[1]{\textbf{#1}\,\textbar\;}

% Figure pages. Most legends are longer than the space a full-width figure leaves
% on one page. Inside a float the surplus is set below the bottom margin and does
% not appear in the output at all, so each figure is placed on a page of its own
% and its legend is ordinary text, free to continue onto the next page. The
% optional argument is the graphic width, \textwidth unless given.
% The graphic is capped in height as well as width so that a tall figure cannot
% push its own legend onto the next page, which leaves a blank one behind it.
\newcommand{\figpage}[3][\textwidth]{%
  \clearpage
  \noindent\makebox[\linewidth]{%
    \includegraphics[width=#1,height=0.82\textheight,keepaspectratio]{#2}}\par
  \vspace{5pt}%
  \begingroup\footnotesize\setlength{\parindent}{0pt}%
  \noindent #3\par\endgroup}

% Four figures are drawn as two artwork files that carry one number and one
% legend between them. They are stacked on a single page and share the legend
% below, which is how they were composed; the second file continues the panel
% lettering of the first.
\newcommand{\figpagetwo}[4][\textwidth]{%
  \clearpage
  \noindent\makebox[\linewidth]{\includegraphics[width=#1]{#2}}\par
  \vspace{3mm}%
  \noindent\makebox[\linewidth]{\includegraphics[width=#1]{#3}}\par
  \vspace{5pt}%
  \begingroup\footnotesize\setlength{\parindent}{0pt}%
  \noindent #4\par\endgroup}

\pagestyle{fancy}\fancyhf{}
\renewcommand{\headrulewidth}{0pt}
\fancyfoot[C]{\sffont\footnotesize\color{muted}\thepage}
